I am writing to express my strong interest in the Los Alamos National Laboratory Quantum Computing Summer School.

\begin{point}{Current research}
	I am currently pursuing a PhD at DESY in Hamburg, where my research focuses on leveraging quantum machine learning techniques for high-energy physics simulations.
	My long-term goal is to develop more efficient and scalable methods for simulating next-generation calorimeters, addressing the exponential growth in data complexity and volume anticipated in future experiments.
\end{point}

\begin{point}{About myself}
	I have built a solid foundation in quantum information science, beginning with my Master's degree, where I focused on Quantum Information, Tensor Networks, and related topics. My interest in these areas led me to further develop my skills through an internship at CERN OpenLab, where I worked on both supervised and unsupervised quantum machine learning models for phase detection in spin systems \cite{ANNNIcern}. In my Master’s thesis, I compared QML methods to Tensor Network approaches for \emph{jet tagging} at LHCb. Currently, as a PhD student at DESY, I am exploring the potential of quantum generative models to simulate high-energy physics events.\\
	In addition to my academic work, my strongest strength in the field lies in my passion for programming in general. I am dedicated to continuously improving my skills. Specifically, in relation to my work, this passion has driven me to develop strong proficiency with many quantum computing libraries such as Pennylane, Qiskit, and Quimb, and I am currently expanding my expertise into Julia libraries.
\end{point}

\begin{point}{Research goals and interests}
	I am particularly interested in rethinking algorithms to handle a larger number of features, which would enable the application of my models to image data. A key challenge to such scaling, especially in generative modeling, is the Barren Plateau problem.\\
	I am indeed particularly drawn to the opportunity to attend the summer school at LANL, as it will allow me to engage with its group, whose work I have been following closely, primarily on the topics of the Barren Plateau and Classical Simulability.
	Particularly, I recently reviewed a paper \cite{zoe} coauthored by Marco Cerezo, which introduced techniques that allowed me to upscale my generative model and embed MNIST images. While I successfully implemented these techniques and developed a training routine, the results have not been satisfactory, likely due to critical trainability issues arising from the large qubit counts I am simulating. I am eager to investigate these challenges further and explore potential solutions.
\end{point}

I am excited about the opportunity to collaborate with the leading team of researchers at LANL, whose contributions have significantly influenced my research path.

Thank you for your consideration.
